%vi: tw=78
\documentclass{article}

%\usepackage{a4wide}
\usepackage[utf8x]{inputenc}
\usepackage[czech]{babel}
\usepackage[IL2]{fontenc}

\pagestyle{empty}

\begin{document}
\begin{flushright}
\today
\end{flushright}
\begin{center}
\Large\textbf{Posudek oponenta na bakalářskou práci \\
\normalsize%
\textit{Michal Krajňanský:} Podpora preprocesoru CPP v nástroji Stanse}
\end{center}

Dle zadání měla bakalářská práce upravit preprocesor jazyka C tak, aby poskytoval mapování preprocesovaného kódu do kódu původního a této informace využít v nástroji Stanse za použití rozhraní, které autor sám navhrne.

Text práce je dobře strukturovaný do 5 kapitol. V úvodu se čtenář dozví
motivaci a k čemu má práce sloužit. Ve druhé kapitole je zbytečně podrobný
popis toho, co preprocesory provádějí, přičemž se tato informace nadále v
textu takřka nevyužije. Tato kapitola nebyla stěžejní a stačilo pouze uvést
odkaz na další zdroje, kde je toto podrobně popsáno.

Následuje zhodnocení 2 preprocesorů a popis jejich instalace. Proč zde autor
uvádí, že preprocesory mají 30000 řádků a jsou obtížné k pochopení? Oba
v práci uvedené preprocesory mají shodně něco přes 16000 řádků a k účelu této
práce se bylo možné omezit jen na malé úseky kódu, které byly předem vytipovány
předchozí diplomovou prací. Komplexita kódu je také podkládána tvrzením, že
\textsc{MCPP} obsahuje \texttt{goto-label} konstrukce. Z mého pohledu je to
běžná praxe, jak ošetřovat chyby na koncích funkcí a na kvalitě, zejm.
čitelnosti, kódu to spíše přidavá.  Fakt, že \textsc{CPP} z \textsc{GCC} je
nejrozšířenější, není taktéž ničím podložen. Co například preprocesor v MS VS?

V této kapitole je také opět věnováno příliš prostoru irelevantním věcem
(rozbalení archívu, kusy kódu interních strukturu atd.). Je zde také zopakován
popis řešení Miloslava Vávry, kde je 1 strana zabrána jeho XSD schématem.

Celá práce je ukončena kapitolou, která navrhuje 3 řešení. Tato kapitola měla
být rozvinuta do více než jedné a půl strany, neboť to bylo to stěžejní, čím
se práce měla zabývat. Navíc příliš nerozumím autorově argumentaci, že jsou
navržené varianty příliš složité a že je třeba rozšiřovat nějaké vnitřní
struktury. Překladač \textsc{GCC} žádné jiné informace nemá a je schopen
ukazovat zpět do zdrojových kódu jen na základě výstupu preprocesoru.

Po jazykové stránce obsahuje práce jen několik překlepů, typografické chyby
(,,s'' a ,,k'' na koncích řádků apod.) a také jazykové tvary ,,debuging'' a
,,kompilace'', pro něž existují české ekvivalenty.

Celkově práce působí nepropracovaným dojmem a klade důraz na nepodstatné.
Nepřináší detailní popis mapování kódu a návrh API pro \textsc{Stanse}, tedy
\textbf{nesplňuje zadání}. Dále, odpočteme-li zbytečný vložený kód, dostaneme
ve výsledku 18 vysázených stran, čímž rozsahem \textbf{nesplňuje požadavky na
bakalářskou práci}. Byla vůbec práce konzultována s vedoucím?

Navrhuji práci \textbf{neuznat} jako bakalářskou, neboť autor
nesplnil podmínky zadání a doporučuji hodnocení stupněm F.

\begin{flushright}
Jiří Slabý
\end{flushright}

\end{document}
